\section{Conclusi\'on}
El uso de \'epocas offline en O.N.O. modula el balance b\'asico entre demandas de comunicaci\'on y fidelidad algor\'itmica del modelo. En la infraestructura local simulada, incrementar $E$ inestabiliz\'o el entrenamiento de LeNet-5 de manera predecible sin aportar ganancias de velocidad, debido a las limitantes de hardware de la CPU hospedadora de 4 cores f\'isicos. 

El estudio evidencia el valor del sistema para diagnosticar fallas de escalabilidad l\'ogica y comportamientos asint\'oticos de optimizadores. Como trabajo futuro, se plantea validar el sistema bajo topolog\'ias multi-nodo reales con latencias de red inducidas por software.
